% page 191 du fac-simile (image p-190 du scan)
% bloc 167.6 x 259.3 mm ; marge gauche 34.4 mm
\pgc{12.700mm}{0.000mm}{
\l{}
\l{}
\l{}
\l{\cel{53}183}
\l{}
\l{fundo. I. La loko maxim infra en to quo esas kava. - II. (\sou{met}af.)}
\l{\cel{3}La parto interna, nevidebla, opoze a lo videbla extera. -III.}
\l{\cel{3}To sur quo ulo su apogas. - IV. La loko maxim fora (en lo}
\l{\cel{3}konsiderata.) - FIS.}
\l{}
\l{fundamento. I. Masonuro quan onu establisas, en sulo, kom bazo}
\l{\cel{3}solida di edifico. - II. (\sou{metaf.}) En ensemblo, la elemento}
\l{\cel{3}esencala sur qua su apogas lo cetera. - To sur quo onu apogas}
\l{\cel{3}sua judiko. - DEFIRS.}
\l{}
\l{\sou{funelo}. I. Implemento ek tubo an qua esas muntita la somito di la}
\l{\cel{3}kono qua formacas lua parto supra, per qua onu varsas liquido}
\l{\cel{3}aden barelo, e c. - II. La parto extera di la orelo. - E.}
\l{}
\l{funero. Ensemblo de la ceremonii qui relatas la sepulto di homo}
\l{\cel{3}jus-mortinta. - DEFIS.}
\l{}
\l{\sou{funesta}. Qua adportas desfortuno. - FIS.}
\l{}
\l{\sou{fungo}. (\sou{bot.}) Planto kriptogama qua, en medio apta, kreskas e}
\l{\cel{3}su reproduktas rapide, e di qua ula speci esas manjebla, la}
\l{\cel{3}ceteri esante venenoza. - EIL.}
\l{}
\l{funikulara. I. (\sou{aludante relvoyo}).Di qua la vagoni, por acensar}
\l{\cel{3}rampo grand-angula, movesas per kordego spulita sur vindilo.}
\l{\cel{3}- II. (\sou{matem.}) Qualeso di la kurvo quan produktas kordo flexe-}
\l{\cel{3}bla, ne-extensebla, suspendita per lua extremaji a du punti}
\l{\cel{3}fixa ed abandonita a la efiko di la gravito ; qualeso di la}
\l{\cel{3}poligono quan formacas kordo sur la punti diversa di qua aplik-}
\l{\cel{3}esas forci. - DEFIS.}
\l{}
\l{\sou{fuorto}. (\sou{milit.}) Laboruro de tero o de masonuro por protektar}
\l{\cel{3}urbo, paseyo, defileo, portuo, e c. - DEFIRS.}
\l{}
\l{\sou{furo}. Pelo de ula animali, kun lua pili, preparita e fasonita}
\l{\cel{3}por garnisar od uzesas kom futero di korpo-vesti, kapo-vesti,}
\l{\cel{3}mufi, e c. - EFS.}
\l{}
\l{\sou{fureto.}\cel{1}(\sou{zool.}) Mikra mamifero karnivora, de la genero \textquotedbl{}martri\textquotedbl{},}
\l{\cel{3}quan onu povas dresar por chasar la kunikli, en la tertruo}
\l{\cel{3}di qua onu enirigas lu. - EFI.}
\l{}
\l{furgono. (\sou{armeo}). Longa veturo, +kluza (klozita), per qua onu}
\l{\cel{3}transportas la bagaji, la provizuri de nutrivi. - F.}
\l{}
\l{\sou{furiar}. (\sou{netrans}.) Kaptesar da iraco dum qua onu ne plus dominacas}
\l{\cel{3}su. - (\sou{metaf.}) Esar violenta extraordinare. - DEFIRS.}
\l{}
\l{\sou{furiero}. I. Oficiro qua, dum voyajo, preiras princo ed agas lo}
\l{\cel{3}necesa pri la lojeyi. - II. (\sou{armeo}). Sub-oficiro qua su okupas}
\l{\cel{3}prokurar la lojeyi necesa por la armeani cirkulanta, repartisar}
\l{\cel{3}la nutrivi, e c. - DFIS.}
\l{}
\l{furnelo. Aparato en qua onu facas fairo, por koquar la nutrivi ;}
\l{\cel{3}ulakaze lu kontenas forno. - FIS.}
}
